% ============================================================================
\section{Discussion}
\label{sec:discussion}

\subsection{What the evidence says, end to end}

\begin{enumerate}
\item \textbf{Temple alone cannot detect seizures.} It measures scalp optical signals over a non-brain artery.
\item \textbf{Seizure chemistry can be simulated and, in principle, inferred.} The ionic model reproduces the known ionic and haemodynamic hallmarks and is validated against Cressman 2009; inference formulas work only when their assumptions hold, and model-based inference beats single-formula conversions. Our Epileptor matches TVB.
\item \textbf{Ear-level EEG sees most seizures but raises too many false alarms on its own.} On CHB-MIT, two ear-position channels matched an 18-channel cap for new users (88\% vs 85\%) with $\approx$2$\times$ the false alarms; per-person normalisation and personal fine-tuning each helped substantially.
\item \textbf{The heart is the body signal that pays.} On SeizeIT2, ECG cut false alarms by 18--44\% at matched sensitivity for 115/125 patients; simulated in-ear PPG keeps most of that benefit.
\item \textbf{Movement matters through the heart, and minutes matter more than seconds.} Heart rate not explained by movement plus 15\,s--5\,min context reduced false alarms by up to 62\% vs ear EEG alone.
\item \textbf{False alarms are personal.} A quarter of users produce most of them; the one linked trait is unstable ear-EEG amplitude (electrode fit); they are predictable from the first days of wear; a personal alarm verifier that learns each user's false alarms is the single most effective strict-budget tool.
\item \textbf{The dangerous seizures are caught.} All focal-to-bilateral tonic-clonic seizures were caught at $\le$5 \FAd.
\item \textbf{The ceiling is the signal.} Short focal seizures ($<$30\,s, over half of all) are nearly invisible to two mastoid channels; a quarter of seizures never raise the detector above 0.5. Better alarm logic cannot close this gap; hardware (contact quality, more or better-placed electrodes) must.
\end{enumerate}

\subsection{Implications for the device (``Prototype 1'')}

The evidence defines a concrete device architecture: bilateral ear-EEG, in-ear PPG, a head IMU and a phone; features computed over a rolling few-minute history; a general model at first use; personal fine-tuning after about two confirmed seizures; a personal alarm verifier; responder-dependent modes (alarm mode for heart responders, with 39\% of seizures caught at $\le$1 \FAd{} and 64\% at $\le$5; a log-only seizure diary for non-responders); a signal-quality gate with a ``reseat'' prompt; user cancellation of alarms (a phase-3 ECG study cancelled all its false alarms this way); and SUDEP-relevant safety alarms (tonic-clonic seizures, bradycardia, desaturation).

Arguments \emph{against} building it now as an alarm: overall performance at $\le$1 \FAd{} (19\%) is not good enough for most users; all numbers come from SeizeIT2's stick-on mastoid electrodes, chest ECG and hospital conditions, while dry ear electrodes in daily life will be worse; electrode fit, the main false-alarm driver, is unmeasured for our own hardware; PPG results are simulated; and a seizure alarm is a regulated medical device whose misses create false reassurance. Arguments \emph{for} building it as a research recording and diary device: only hardware can answer the largest open questions (dry-electrode signal quality, real in-ear PPG while walking, chewing and sleeping); tonic-clonic seizures are already caught; objective seizure counts are clinically valuable because patients under-report; all components exist off the shelf and the software pipeline is ready. The recommendation is to build it first as a recording and seizure-diary device, validated in phases (bench, self-wear by healthy volunteers, then an epilepsy-monitoring-unit partnership with ethics approval), and to enable alarms only when device-specific data support them.

\subsection{Limitations}

\begin{itemize}
\item \textbf{Simulation:} ictal events are not self-sustaining at the default pump setting; the network is small, 1-D, single-type per population, with simplified conservation, implicit energy and generic haemodynamics; the LFP is a proxy.
\item \textbf{Inference:} all results are simulation-only; T2 is an inverse crime; E2's DC channel needs a calibrated gain.
\item \textbf{CHB-MIT:} paediatric, hospital, mostly at rest; ear montages were approximated with scalp channels (T7--P7, T8--P8), not dry ear electrodes; small patient-specific samples.
\item \textbf{SeizeIT2:} hospital recordings (less movement than daily life); chest ECG, not ear PPG (PPG simulated at feature level); mastoid electrodes with adhesive, not a hook; best-along-the-curve numbers slightly optimistic; fine-tuned models need held-out calibration; personalised experiments cover 71 patients with $\ge$3 seizures; responder status needs ECG and two seizures; age unavailable.
\item \textbf{Live Lab:} body signals follow literature timings, not the neuron model; the fusion rule and \Kp-DC threshold are hand-set.
\item \textbf{Not a medical device.}
\end{itemize}

\subsection{Future work}

(1) Finish the raw-EEG CNN and stacking experiment. (2) Personal verifier on top of the fine-tuned detector with calibration on a held-out slice of the tuning period. (3) Hardware prototype: dry ear electrodes with contact-quality sensing, in-ear PPG/SpO\textsubscript{2}, IMU, phone app with cancellation. (4) Real in-ear PPG recordings to replace the simulation. (5) Sequence models over minutes of HR/HRV and EEG (pre-ictal HRV change occurs in $\approx$41\% of seizures). (6) Seizure forecasting from multi-day cycles. (7) Simulation: recalibrate the pump current for self-sustained seizures, add ATP/O\textsubscript{2}, interneuron classes, a 2-D sheet and a volume-conductor EEG forward model. (8) Inference on real data: fit the ionic model to patients with simultaneous intracranial and scalp recordings.

\section{Conclusion}

Starting from a question about a temple-worn optical wearable, this project built the scientific and engineering basis for a different device: a bilateral ear-hook combining ear-EEG, in-ear PPG and motion, with a personalised, multimodal, two-stage detector. On the largest public wearable dataset, the heart signal, minute-scale context, a movement-adjusted heart-rate feature and personalisation together raised the share of seizures caught at $\le$1 false alarm/day from 4\% (ear EEG) to 19\% overall and 39--40\% in heart responders, caught every focal-to-bilateral tonic-clonic seizure at $\le$5 false alarms/day, and remained robust to simulated PPG. The remaining limitation is the information available from two behind-the-ear electrodes, which points the next phase toward hardware.
